\section{Задача и данные}

\subsection{Постановка}
По текстовому условию нужно сгенерировать Python-код с функцией
\code{solve()}, возвращающей ответ. По правилам задания модель обучается
с нуля и имеет Transformer-энкодер и декодер. Все ответы официального
теста должны генерироваться за 180 секунд на одной NVIDIA T4;
дополнительные обучающие данные и таргеты запрещены~\cite{assignment}.
Также я учитывал отдельный запрет спекулятивного декодирования.

Цель экспериментов --- выбрать конфигурацию, которая сочетает качество
решений, устойчивое обучение и приемлемую стоимость генерации.
Названия BART и T5 далее обозначают устройство моделей:
предобученные веса не использовались.

\subsection{Подготовка и разбиения}
Использован корпус \code{SlavaYus/tinysms_qa_12m}: в исходном train
11\,838\,003 пары \code{question}/\code{code}, официальный test содержит
1919 условий без кода~\cite{dataset}. Комментарии в коде удалялись
Python-токенизатором с сохранением строковых литералов и структуры
программы. Длинные обучающие пары исключались целиком, а не обрезались.

Первые эксперименты использовали разбиение 80/10/10 с seed 42.
Начиная с модели 002 использовалось 98/1/1. После фильтрации по лимитам
256 токенов осталось 11\,593\,060 train и 118\,302 validation.
Новый внутренний test выделен отдельно из исходного train.
Официальный test при обучении и подборе параметров не использовался.

\begin{table}[htbp]
\centering\small
\begin{tabularx}{\linewidth}{lXr}
\toprule
Обозначение & Назначение & Размер \\
\midrule
Test-001 & Выборка из test разбиения 80/10/10 & 10\,000 \\
Test-002 & Выборка из нового test разбиения 98/1/1 & 10\,000 \\
Val-20k & Фиксированная validation коротких сравнений & 20\,000 \\
Val-full & Полная validation моделей 002/003 и нового BART & 118\,302 \\
Official & Официальный test без открытых эталонов & 1919 \\
\bottomrule
\end{tabularx}
\caption{Наборы с разными назначениями. Test-001 и Test-002 нельзя считать
одной выборкой, хотя оба содержат 10\,000 строк.}
\label{tab:splits}
\end{table}

Для короткого поиска взяты фиксированные 1\,000\,000 train и Val-20k.
В последующем сравнении данных использовались повторение этого миллиона
четыре раза либо четыре миллиона различных исходных train-строк.
Полный новый BART использовал прежний токенизированный кэш целиком:
порядок строк и исходные таргеты сохранены, ретокенизация не выполнялась.

Разбиение выполнено по строкам. Группировка одинаковых условий между
частями не реализована, поэтому возможны повторы и зависимые примеры.
Отдельно учтено, что старый Test-001 после перехода к 98/1/1 уже
не является гарантированно отложенной выборкой для модели 002.

\subsection{Токенизатор}
Обучен собственный byte-level BPE со словарём 8000 токенов на разрешённых
train-текстах; готовый словарь другой модели не загружался.
Байтовое представление сохраняет отступы и переводы строк программы.
Спецтокены: PAD, UNK, BOS и EOS с ID 0, 1, 2, 3.
В labels нет начального BOS, padding заменён на $-100$;
декодер начинает генерацию с BOS и заканчивает по EOS либо лимиту длины.
В сравнениях архитектур использован один и тот же сохранённый BPE.

\section{Модели, обучение и измерения}

\subsection{Архитектуры}
Основная реализация использует BART encoder--decoder из Transformers,
инициализированный случайно~\cite{bart}. Энкодер получает условие,
декодер --- предыдущие токены программы с причинной маской.
Embeddings энкодера, декодера и выходной проекции связаны.
Обозначение E+D далее задаёт число слоёв энкодера и декодера.

\begin{table}[htbp]
\centering\small
\begin{tabularx}{\linewidth}{lrrrX}
\toprule
Модель & E+D & $d$ & FFN & Особенности \\
\midrule
Baseline-001 & 6+6 & 512 & 2048 & BART, 8 heads \\
Transformer-002 & 8+8 & 768 & 3072 & 139.25 млн параметров \\
Transformer-003 & 10+10 & 1024 & 4096 & 303.18 млн параметров \\
BART-кандидат & 8+8 & 1024 & 4096 & 244.39 млн параметров \\
T5-style & 8+8 & 1024 & 2816 & GEGLU, Pre-RMSNorm, relative bias \\
Copy & 8+8 & 1024 & 4096 & BART и отдельная copy attention \\
\bottomrule
\end{tabularx}
\caption{Основные архитектуры. Различия глубины в коротком поиске
дополнительно приведены в таблице~\ref{tab:h100}.}
\end{table}

T5-style использует относительный attention bias с 32 buckets,
GEGLU и связанные embeddings; число параметров --- 247\,310\,848.
Вариант вдохновлён семейством T5~\cite{t5}, но не воспроизводит полностью
рецепт T5 v1.1: здесь сохранены dropout 0.1 и tied embeddings.
У BART attention dropout равен 0.0, поэтому сравнение затрагивает
весь набор архитектурных настроек, а не одну нормализацию.

В Copy к BART добавлено распределение копирования по токенам условия.
Идея pointer-generator~\cite{pointer} реализована как смесь
\[
 p(y_t\mid x,y_{<t})=g_t\,p_{\mathrm{vocab}}(y_t)
 +(1-g_t)\sum_{i:x_i=y_t}a_{ti}.
\]
Повторяющиеся токены входа суммируются, специальные и masked токены
исключаются. При отсутствии допустимых токенов используется словарь.
Дополнительные метки копирования не нужны; обучение использует прежний код.
Число параметров Copy --- 244\,656\,129.

